\chapter{Details}\label{ch:Details}

\section{Werkzeuge}\label{sec:Werkzeuge}
\begin{mydefinition}[$bin$ und $dec$]\label{mydefinition:binunddec}
\leavevmode
\begin{itemize}
	\item Es sei $n>0$ eine natürliche Zahl. Dann bezeichnen wir mit $bin(n)$ die binäre (Basis 2) Zahlendarstellung von $n$ ohne führende Nullen.
	\item Es sei $n>0$ eine natürliche Zahl. Dann bezeichnen wir mit $dec(n)$ die dezimale (Basis 10) Zahlendarstellung von $n$ ohne führende Nullen.
	\item Wir definieren zusätzlich $bin(0) := 0$ sowie $dec(0) := 0$.
\end{itemize}
\end{mydefinition}


\begin{mytheorem}[geometrische Summe]\label{mytheorem:geometrischeSumme}
	Es sei $q \neq 0$ eine reelle Zahl und $n$ eine natürliche Zahl. Dann gilt
	\begin{align*}
		\sum_{k = 0}^n q^k = q^0 + q^1 + \ldots + q^n = \frac{q^{n+1}-1}{q-1}.
	\end{align*}
	\begin{myproof}
		Wir beweisen die Aussage durch vollständige Induktion.
		\begin{itemize}
			\item Für $n = 0$ gilt die Aussage, da
			      \begin{align*}
				      \sum_{k = 0}^0 q^k = q^0 = 1 = \frac{q^{0+1}-1}{q-1}.
			      \end{align*}
			\item Die Aussage gelte nun für eine natürliche Zahl $n$. Wir zeigen, dass sie auch für $n+1$ gilt.
			      \begin{align*}
				      \sum_{k = 0}^{n+1} q^k = q^{n+1} + \sum_{k = 0}^n q^k = q^{n+1} + \frac{q^{n+1}-1}{q-1} = \frac{q^{n+2}-1}{q-1}.
			      \end{align*}
		\end{itemize}
	\end{myproof}
\end{mytheorem}


\begin{mytheorem}[verallgemeinerte geometrische Summe]\label{mytheorem:verallgemeinertegeometrischeSumme}
	Es sei $q \neq 0$ eine reelle Zahl und $n$ und $m\leq n$ natürliche Zahlen. Dann gilt
	\begin{align*}
		\sum_{k = m}^n q^k = q^m + q^{m+1} + \ldots + q^n = \frac{q^{n+1}-q^m}{q-1}.
	\end{align*}
	\begin{myproof}
		Unter Verwendung von \cref{mytheorem:geometrischeSumme} finden wir
		\begin{align*}
			\sum_{k = m}^n q^k =
			\sum_{k = 0}^n q^k - \sum_{k = 0}^{m-1} q^k =
			\frac{q^{n+1}-1}{q-1} - \frac{q^m - 1}{q-1} =
			\frac{q^{n+1}-q^m}{q-1}.
		\end{align*}
	\end{myproof}
\end{mytheorem}


\begin{mytheorem}[Anzahl Ziffern in Zahlendarstellung]\label{mytheorem:AnzahlZiffern}
	Es seien $n>0$ und $b>1$ natürliche Zahlen. Die kürzeste $b$-adische Darstellung von $n$ (Darstellung ohne führende Nullen) hat genau $\ceil{\log_b(n + 1)}$ Stellen.
	\begin{myproof}
		Es sei $s$ die Anzahl der Stellen der kürzesten $b$-adischen Darstellung von $n$. Die grösste $s$-stellige Zahl in Basis $b$ (kürzeste Darstellung) ist
		\begin{align*}
			 & \sum_{k = 0}^{s-1} (b-1)b^{k} = (b-1)\sum_{k = 0}^{s-1} b^{k} = (b-1)\frac{b^s - 1}{b - 1} = b^s - 1,
		\end{align*}
		wobei wir \cref{mytheorem:geometrischeSumme} verwendet haben. Die kleinste $s$-stellige Zahl Basis $b$ (kürzeste Darstellung) ist $b^{s-1}$. Somit gilt
		\begin{align*}
			b^{s-1} - 1 < n \leq b^s - 1 \iff \\
			b^{s-1} < n + 1 \leq b^s \iff     \\
			s-1 < \log_b(n+1) \leq s,
		\end{align*}
		wobei wir verwendet haben, dass $\log_b$ eine streng monoton wachsende Funktion ist.
		Dann folgt aber $\ceil{\log_b(n + 1)} = s$.
	\end{myproof}
\end{mytheorem}



\begin{mytheorem}[Obere Schranke binäre Darstellungslänge]\label{mytheorem:Estimate}
	Für jede natürliche Zahl $n>0$ gilt
	\begin{align*}
		\ceil{\log_2(n+1)}\leq \ceil{\log_2(n)} + 1.
	\end{align*}
	Offensichtlich folgt aus dieser Behauptung sofort
	\begin{align*}
		\ceil{\log_2(n+1)}\leq \ceil{\log_2(n)} + 1 < \log_2(n) + 2.
	\end{align*}
	\begin{myproof}
		Wir beweisen zunächst die Ungleichung $\log_2(n+1) \leq \log_2(n) + 1$. Dazu berechnen wir zunächst
		\begin{align*}
			\log_2(n) + 1 = \log_2(n) + \log_2(2) = \log_2(2n).
		\end{align*}
		Damit gilt also
		\begin{align*}
			\log_2(n+1) \leq \log_2(n) + 1 & \iff \\
			\log_2(n+1) \leq \log_2(2n)    & \iff \\
			n + 1\leq 2n                   & \iff \\
			1 \leq n,
		\end{align*}
		wobei wir verwendet haben, dass $\log_2$ (streng) monoton wachsend ist. Die Ungleichung ist für alle natürlichen Zahlen $n>0$ korrekt.
		Wir berechnen nun
		\begin{align*}
			\log_2(n+1) \leq \log_2(n) + 1               & \Rightarrow             \\
			\ceil{\log_2(n+1)} \leq \ceil{\log_2(n) + 1} & = \ceil{\log_2(n)} + 1.
		\end{align*}
	\end{myproof}
\end{mytheorem}

\begin{mydefinition}
	Es sei $\Sigma = \lrc{a_1, a_2 , \ldots, a_n}$, $n > 0$, ein Alphabet mit der Ordnung $a_1 < a_2 < \ldots < a_n$. Es seien $x, y$ beliebige Wörter über $\Sigma$. Wir definieren die kanonische Ordnung $<$ auf allen Wörter über $\Sigma$ wie folgt:
	\begin{enumerate}
		\item Falls $x$ kürzer ist als $y$, dann gilt $x < y$.
		\item Sind anderenfalls $x$ und $y$ gleich lang, dann gilt $x < y$ genau dann, wenn $x$ alphabetisch vor $y$ liegt.
	\end{enumerate}
\end{mydefinition}

\section{Code}

\lstinputlisting[language=Python,caption=\texttt{mandelbrot.py},label=listing:mandelbrot]{EF/Kolmogorov/mandelbrot.py}
